% Release notes: pdf-to-markdown v2.5
% Compile with: xelatex pdf-to-markdown_v2_5_release_notes.tex  (run twice)
\documentclass[11pt,a4paper]{article}

\newcommand\skillname{pdf-to-markdown}
\newcommand\skillversion{2.5}
\newcommand\skilldate{12 October 2026}

% ---------------------------------------------------------------- layout
\usepackage[margin=2.4cm,top=2.7cm,bottom=2.6cm,headheight=14pt]{geometry}
\usepackage{amssymb}
\usepackage{fontspec}
\setmainfont{Latin Modern Roman}
\setsansfont{Latin Modern Sans}
\setmonofont{DejaVu Sans Mono}[Scale=0.82]
\usepackage{microtype}
\usepackage{xcolor}
\definecolor{accent}{HTML}{1F5F7A}
\definecolor{codebg}{HTML}{F2F5F7}
\definecolor{muted}{HTML}{6B7780}
\usepackage{titlesec}
\titleformat{\section}{\sffamily\Large\bfseries\color{accent}}{\thesection}{0.8em}{}
\titleformat{\subsection}{\sffamily\large\bfseries}{\thesubsection}{0.7em}{}
\titlespacing*{\section}{0pt}{1.5em}{0.5em}
\titlespacing*{\subsection}{0pt}{1.1em}{0.35em}
\usepackage{enumitem}
\setlist{itemsep=0.2em, topsep=0.35em}
\usepackage{booktabs,tabularx,array}
\renewcommand\tabularxcolumn[1]{>{\raggedright\arraybackslash}p{#1}}
\setlength\extrarowheight{1pt}
\usepackage[most]{tcolorbox}
\tcbuselibrary{listings}
\newtcblisting{cmd}{listing only, colback=codebg, colframe=codebg, boxrule=0pt,
  arc=2pt, left=6pt, right=6pt, top=2pt, bottom=2pt, before skip=6pt, after skip=8pt,
  listing options={basicstyle=\ttfamily\small, breaklines=true, columns=fullflexible,
  keepspaces=true, upquote=true, breakindent=1.5em}}
\newtcolorbox{glance}{enhanced, colback=accent!5, colframe=accent, boxrule=0.6pt,
  arc=2pt, left=9pt, right=9pt, top=6pt, bottom=6pt,
  fonttitle=\sffamily\bfseries, title=At a glance}
\newtcolorbox{note}{enhanced, colback=accent!4, colframe=accent, boxrule=0pt,
  leftrule=2.5pt, arc=0pt, left=8pt, right=8pt, top=5pt, bottom=5pt,
  before skip=8pt, after skip=8pt}
\usepackage{fancyhdr}
\pagestyle{fancy}
\fancyhf{}
\fancyhead[L]{\sffamily\small\color{muted}\skillname\ v\skillversion}
\fancyhead[R]{\sffamily\small\color{muted}Release notes}
\fancyfoot[C]{\sffamily\small\color{muted}\thepage}
\renewcommand\headrulewidth{0.3pt}
\usepackage[colorlinks=true, linkcolor=accent, urlcolor=accent]{hyperref}
\hypersetup{pdftitle={\skillname\ v\skillversion\ release notes},
  pdfsubject={Claude skill release notes}}
\DeclareUrlCommand\code{\urlstyle{tt}}
\newcommand\step[1]{\textsf{\textbf{#1}}}

\begin{document}

% ---------------------------------------------------------------- title
\thispagestyle{plain}
\begin{flushleft}
{\sffamily\small\color{muted}\MakeUppercase{Claude skill \textbullet{} Release notes}}\\[4pt]
{\sffamily\fontsize{26}{30}\selectfont\bfseries\color{accent}\skillname}\\[6pt]
{\sffamily\large Version \skillversion\quad\textcolor{muted}{|}\quad \skilldate}\\[2pt]
{\color{accent}\rule{\linewidth}{0.8pt}}
\end{flushleft}

\begin{glance}
\begin{tabularx}{\linewidth}{@{}l X@{}}
\textbf{What it is} & A converter that turns a PDF of a paper, lecture notes or textbook
  chapter into text-only Markdown with exact LaTeX mathematics, plus one companion PDF of its
  figures. Every step is machine-checked. \\
\textbf{When it runs} & When you ask Claude to convert a PDF to Markdown (\code{.md}). \\
\textbf{What you get} & \code{NAME.md} and \code{NAME_figures.pdf}, ready to upload to a
  Project. \\
\textbf{Requires} & Claude with skills and code execution enabled. PyMuPDF and KaTeX are
  installed on first use; pdflatex is used when present. \\
\textbf{Install} & \textsf{Customize \textrightarrow{} Skills} (Section~\ref{sec:install}).
\end{tabularx}
\end{glance}

% ---------------------------------------------------------------- 0
\section{What's new in 2.5}\label{sec:new}

Version 2.5 is a maintenance release. It fixes seven defects in the helper scripts that
surfaced when converting a 21-page two-column APS paper with 131 numbered equations and six
figures. The conversion procedure, the output format and the writing conventions are
unchanged; transcriptions made with 2.4 remain valid.

\begin{center}
\small
\begin{tabularx}{\linewidth}{@{}l X@{}}
\toprule
\textbf{Script} & \textbf{Fix} \\
\midrule
\code{prepare.py} & \textbf{Printed page numbers decoded.} Page labels stored as UTF-16
  strings (a byte-order mark followed by the text) were reported undecoded, e.g.\
  \code{<FEFF0031>} for page~1. They are now decoded, and a document whose labels simply
  repeat the PDF page numbers is reported as such. \\
\code{prepare.py} & \textbf{Equation tags found in every column.} Tags were accepted only in
  the right quarter of the page, so every tag in the left column of a two-column layout was
  missed (63 of 131 in the test paper). Text columns are now detected per page, and a tag
  is accepted when it is flush with the right edge of its own column. Inline references
  such as ``Eq.~(7)'' at the end of a prose line are still rejected. Latin Modern math fonts
  (\code{LMMath*}) are now recognised as math fonts. \\
\code{figures.py} & \textbf{Figure crops kept inside their column.} Figures embedded as form
  objects often declare a bounding box far larger than what they draw, and the crop spilled
  into the neighbouring column; page numbers and running heads above a figure were also
  included. The crop is now clipped to the caption's column, excludes page numbers and
  running heads, and is trimmed to the pixels actually drawn. \\
\code{figures.py} & \textbf{No captions from running text.} A caption now needs a separator
  after its number (\emph{FIG.\,1.}, \emph{Figure~1:}, \emph{Fig.\,1~|}), so a sentence
  beginning ``Figure~1(a), computed for\,\ldots'' is no longer taken for a second Fig.~1. \\
\code{check.py} & \textbf{No false failures from missed tags.} A \verb|\tag| whose ``(N)'' is
  in the text layer but was not detected as a tag is listed for a check by eye instead of
  failing the run. With the tag fix above this list is normally empty. \\
\code{verify_eqs.py} & \textbf{Every tagged equation on the sheets.} With all tags detected,
  no equation is left off the proofreading sheets (68 of 131 were before). \\
\code{verify_eqs.py} & \textbf{Proofreading crops show the whole equation.} Crops spanned the
  full text width and were bounded by lines of the other column, so they were about one line
  high and misaligned. A crop now covers only the equation's own column, between the nearest
  prose lines of that column (an equation row that starts at the margin is not mistaken for
  prose), with room for multi-line displays and over-wide rows. \\
\bottomrule
\end{tabularx}
\end{center}

\noindent Effect on the test paper:
\begin{center}
\small
\begin{tabular}{@{}l l l@{}}
\toprule
 & \textbf{2.4} & \textbf{2.5} \\
\midrule
Equation tags detected & 68 of 131 & 131 of 131 \\
\code{check.py} & FAIL (63 false misses) & PASS \\
Equations on proofreading sheets & 68 & 131 \\
Figure crops correct without manual rework & 2 of 6 & 6 of 6 \\
False captions & 1 & 0 \\
\bottomrule
\end{tabular}
\end{center}
A single-column test document gives the same tags and figure as in 2.4, plus the corrected
caption handling.

\begin{note}
\textbf{Upgrading.} Replace the installed 2.4 with the 2.5 ZIP (Section~\ref{sec:install}).
\code{SKILL.md} changes only where it describes what the scripts do; there is nothing new to
learn and no need to reconvert existing transcriptions.
\end{note}

% ---------------------------------------------------------------- 1
\section{Overview}

Mathematics extracted from a PDF's text layer is unreliable. Calligraphic, roman and bold
letters merge; accents detach; sub- and superscripts flatten; integral and sum signs turn into
placeholder codes. Reading every equation from page images avoids this, but repeating that
work in every conversation is slow and costs a lot of context.

\skillname{} does that work \emph{once}. Claude transcribes each page from rendered images
into Markdown, writing every formula, inline or displayed, as LaTeX read off the image. The
result is then checked by script and proofread against the original. After conversion,
Claude can read the document as clean source in any later conversation, at a fraction of the
cost of the original PDF and without the risk of a corrupted text layer.

% ---------------------------------------------------------------- 2
\section{Outputs}

Both files are named after the original PDF exactly as uploaded, with the same characters,
case, spaces and underscores.

\subsection{\texorpdfstring{\texttt{NAME.md}}{NAME.md}: the transcription}
\begin{itemize}
\item \textbf{Text only.} No image data is embedded, so the file stays small and fully searchable.
\item \textbf{A legend on the first line} identifies the file as a transcription and explains its page markers.
\item \textbf{Page markers.} Every PDF page opens with \texttt{<!-{}- page N -{}->}, or
  \texttt{<!-{}- page N (printed M) -{}->} when the number printed on the page differs from its
  position in the file. Answers drawn from the transcription can therefore cite pages. Where a
  paragraph runs across a page break, it is rejoined and the marker stands inline at the break.
\item \textbf{Exact mathematics.} Inline math uses \code{$...$}. Displays use
  \code{$$...$$} with \verb|\tag{N}| carrying the \emph{printed} equation number, even
  when numbers are out of sequence. Typeface is reproduced exactly (\verb|\mathcal|,
  \verb|\mathrm|, \verb|\boldsymbol|, upright constants where printed). Only KaTeX- and
  MathJax-compatible commands are used, with no custom macros.
\item \textbf{Faithful to the original.} Headings, run-in paragraph heads and the document's
  own cross-reference style are kept. Apparent typos are reproduced, not corrected.
\item \textbf{Figure placeholders.} Each figure becomes a line naming where to find it,
  followed by its transcribed caption.
\end{itemize}

\noindent An illustrative excerpt:
\begin{cmd}
<!-- Transcription of example.pdf. "page N" markers give the PDF page on which the following text starts; "printed M" is the page number printed on that page. -->

<!-- page 1 (printed 213) -->

## 2.1 Linear response

The retarded correlation function is

$$ C^{\mathrm{R}}(t) = -\mathrm{i}\,\theta(t)\,\langle [\hat A(t), \hat B(0)] \rangle \tag{2.14} $$

*[Fig. 3: example_figures.pdf, page 1]*

**Fig. 3.** Real and imaginary parts of the response function.
\end{cmd}

\subsection{\texorpdfstring{\texttt{NAME\_figures.pdf}}{NAME\_figures.pdf}: the figures}
All figures are packed into one PDF at native size and in order of appearance. Each is
labelled with its figure number and the document page it comes from. Packing is adaptive:
figures are added to a page only while every figure keeps at least 150\,dpi effective
resolution when Claude views the page, so small labels stay legible at the lowest token cost.
A document without figures gets no figures PDF. Text inside figures (axis labels, legends,
diagram labels) is part of the image and is not transcribed. Every caption is in the
Markdown, so the figures PDF can be left out of a Project when no figure needs to be seen.

% ---------------------------------------------------------------- 3
\section{How it works}

\subsection{Ground rules}
\begin{itemize}
\item \textbf{No mathematics from the text layer.} Every formula is transcribed from the page image.
\item \textbf{The text layer is a lookup, not a source.} It is used only to place strip cuts,
  find tags, page numbers and captions, and resolve a single doubtful prose word.
\item \textbf{Mandatory zooms} for the distinctions most often misread: bars spanning several
  characters, tildes and hats, primes, the extent of under- and overbraces (including a
  leading minus sign), bold versus regular, and calligraphic versus roman versus italic.
\item \textbf{Every token is written once.} Pages go to separate files, the document is joined
  by script, and corrections are targeted edits. The document is never re-output in full.
\item \textbf{Source first.} If the LaTeX source is available, Claude converts it with pandoc
  and normalises the result instead of transcribing images.
\end{itemize}

\subsection{Pipeline}
\begin{center}
\small
\begin{tabularx}{\linewidth}{@{}l l X@{}}
\toprule
\textbf{Step} & \textbf{Tool} & \textbf{What happens} \\
\midrule
1 Prepare & \code{prepare.py} & Renders each page at 170\,dpi as horizontal strips, cut in the
  gaps between lines so no line or equation is sliced. Reports per page whether it contains
  math, the strips' positions, the printed equation tags (found at the right edge of each
  text column, in single- and multi-column layouts), and the printed page number (from the
  PDF's page labels or a consistent header/footer offset). \\
2 Figures & \code{figures.py} & Locates each figure from its caption (\emph{Fig.\,N},
  \emph{Figure\,N}, \emph{FIG.\,N}, followed by a separator) and renders the whole region,
  vector art and embedded images together, at 200\,dpi, kept inside the caption's column and
  trimmed to what is drawn. Writes a check image with the crop outlined, which Claude
  inspects before going on. \\
3 Transcribe & Claude & Views each page's strips in order and writes that page to its own
  Markdown file. Small or ambiguous symbols are re-rendered at 500--600\,dpi. \\
4 Assemble & \code{assemble.py} & Joins the pages, rejoins paragraphs split by page breaks,
  adds the legend and page markers, packs the figures PDF and resolves figure placeholders. \\
5 Check & \code{check.py} & Runs the mechanical checks (Section~\ref{sec:checks}). Claude
  fixes each finding with a targeted edit. \\
6 Proofread & \code{verify_eqs.py} & Builds proofreading sheets for the displayed equations;
  Claude compares them symbol by symbol. \\
7 Deliver & Claude & Presents \code{NAME.md} and \code{NAME_figures.pdf}, with a short report. \\
\bottomrule
\end{tabularx}
\end{center}

\subsection{Quality checks}\label{sec:checks}
\begin{itemize}
\item \code{check.py} verifies that every formula parses in KaTeX, that every printed equation
  tag appears exactly once, that no text-layer debris (such as \code{(cid:NN)} codes) remains,
  and that every PDF page has exactly one marker, in order. A tag that is printed but was
  not detected as a tag is listed for a check by eye rather than failed.
\item \code{verify_eqs.py} typesets each displayed equation from the transcription with
  pdflatex and places it under the crop of the same equation from the PDF, four per sheet.
  The crop covers the equation's own column, between the nearest prose lines.
  Equations that cannot be placed on a sheet are reported as problems to fix.
\item Inline mathematics is not covered by the sheets. Claude re-reads it on the strips for
  any line with indices or accents.
\item The final report states pages converted (with the printed range), figures and
  figures-PDF pages, check results, and any symbol left ambiguous, with its location.
\end{itemize}

% ---------------------------------------------------------------- 4
\section{Requirements}
\begin{itemize}
\item A Claude plan with skills available and \textbf{code execution} enabled in Claude's settings.
\item Python 3 with \textbf{PyMuPDF} and Pillow. PyMuPDF is installed with pip on first use if missing.
\item \textbf{Node.js} with \textbf{KaTeX}, for the syntax check. KaTeX is installed with npm
  on first use if missing. These installs need the code-execution environment to reach the
  Python and npm package registries.
\item \textbf{pdflatex} for the proofreading sheets. If it is absent, that step is skipped and
  Claude proofreads from the strips instead.
\item \textbf{pandoc}, only when converting from LaTeX source.
\end{itemize}

% ---------------------------------------------------------------- 5
\section{Installation}\label{sec:install}

\begin{enumerate}
\item \step{Package the skill.} The skill is a folder named \code{pdf-to-markdown} containing
  \code{SKILL.md} and \code{scripts/}. If you received it as a ZIP, use that file as it is.
  Otherwise, zip the folder so that the archive contains the folder itself, with
  \code{SKILL.md} directly inside it:
\begin{cmd}
zip -r pdf-to-markdown.zip pdf-to-markdown
\end{cmd}
\item \step{Enable code execution} in Claude's settings if it is not already on. If your
  organisation restricts network access from code execution, make sure the Python and npm
  package registries are allowed.
\item \step{Upload.} Open \textsf{Customize \textrightarrow{} Skills}, choose to upload a
  skill, and select the ZIP.
\item \step{Turn it on.} Check that \emph{pdf-to-markdown} appears in the list and is enabled.
\item \step{Verify.} In a new chat, attach a short PDF with a few equations and ask
  \emph{``Convert this PDF to Markdown.''} Claude should read the skill, run \texttt{prepare.py},
  and finish by presenting \code{NAME.md} (and \code{NAME_figures.pdf} if the document has
  figures).
\end{enumerate}

\begin{note}
\textbf{Updating.} Keep only one copy installed. Remove the previous version and upload the
new ZIP in its place. The installed version is recorded in the \code{metadata.version} field
of \code{SKILL.md}. Transcriptions made with earlier versions remain valid files; reconvert a
document only if you want it in the current output format.
\end{note}

% ---------------------------------------------------------------- 6
\section{Usage}

\subsection{Converting}
Attach the PDF and ask for a conversion, for example:
\begin{itemize}
\item \emph{``Convert this paper to Markdown.''}
\item \emph{``Convert Chapter 3 (pages 41--68) of this book to Markdown.''}
\item \emph{``Here is the PDF and its LaTeX source; convert it to Markdown.''}
\end{itemize}
Download both files when the conversion finishes.

\subsection{Using the result in a Project}
Upload \code{NAME.md}, and \code{NAME_figures.pdf} if you need the figures, to the Project's
knowledge. You can remove the original PDF from the Project or keep it alongside; if both are
present, the transcription is used for content and the PDF only for questions about the page
itself. With the companion skill \textbf{academic-pdf-reading} installed, Claude recognises the
transcription automatically, quotes it as source and cites its page markers.

\subsection{Budget}
Conversion is the expensive step, and it happens once.
\begin{center}
\small
\begin{tabular}{@{}l r@{}}
\toprule
\textbf{Item} & \textbf{Approximate cost} \\
\midrule
One 170\,dpi strip & 600 tokens \\
One page (4--5 strips) & 2{,}500--3{,}000 tokens, plus zooms \\
One proofreading sheet (four equations) & 1{,}500 tokens \\
Figure crop checks & a few hundred tokens each \\
\bottomrule
\end{tabular}
\end{center}
A ten-page note fits comfortably in one conversation. For a long paper or a book, convert one
chapter or section per conversation; Claude says so at the start.

% ---------------------------------------------------------------- 7
\section{Limitations}
\begin{itemize}
\item \textbf{Captions drive figure detection.} Figures without a \emph{Fig.}/\emph{Figure}
  caption, or whose automatic crop is wrong, are rendered by hand, which Claude does when the
  crop check shows a problem.
\item \textbf{Figure text is not transcribed.} Questions about what a figure shows need the figures PDF.
\item \textbf{Inline math is proofread on the strips}, not on the sheets, so it relies on the
  mandatory zooms and a re-read rather than a side-by-side comparison.
\item \textbf{Oversized figures} that fall below 150\,dpi even alone on a page are reported
  with a warning, not shrunk.
\item \textbf{Born-digital PDFs are the intended input.} The scripts use the text layer for
  geometry (strip cuts, tags, page numbers, captions). Scanned documents without a usable text
  layer get less help from them.
\item \textbf{Section-numbered equation tags} such as (2.14) are not detected as tags. The
  transcription still carries them in \verb|\tag{2.14}| and \code{check.py} lists them for a
  check by eye, but they do not reach the proofreading sheets; Claude proofreads them on the
  strips.
\item \textbf{Proofreading crops may include a neighbouring line} when the line just above or
  below an equation is very short (for example ``with $z \in \mathbb{C}$.''). This is harmless.
\item \textbf{Long documents must be split} across conversations because of the context cost.
\item \textbf{Faithful, not corrected.} Typos and inconsistencies in the original are reproduced as printed.
\end{itemize}

% ---------------------------------------------------------------- 8
\section{For maintainers: regression tests}
When a LaTeX source exists for a test document, two scripts measure conversion quality. They
work on any tool's Markdown output, not only this skill's.
\begin{cmd}
python3 scripts/eval/mathcmp.py SOURCE.tex NAME.md [--lenient]
python3 scripts/eval/prosecmp.py SOURCE.tex NAME.md
\end{cmd}
\code{mathcmp.py} compares every display and inline formula token by token, after expanding
zero-argument \verb|\newcommand|s. The \code{--lenient} option also forgives upright versus
italic i, d and e. Count exact matches: mean similarity stays high even when formulas are
wrong. \code{prosecmp.py} compares the prose word by word. Remaining differences are often
spelling-only (\verb|F_{\rm rest}| versus \verb|F_\mathrm{rest}|) and should be judged
individually.

% ---------------------------------------------------------------- 9
\section{Working with academic-pdf-reading}
The two skills are designed as a pair. \textbf{pdf-to-markdown} spends the image-reading effort
once and produces a transcription in a documented format: a legend line, page markers with
printed page numbers, \verb|\tag|ged equations and figure placeholders.
\textbf{academic-pdf-reading} recognises that format, reads it as source without re-verifying
it against page images, opens the figures PDF for figure questions, and escalates only when
the transcription is contradicted or structurally defective.

% ---------------------------------------------------------------- 10
\section{Package contents}
\begin{center}
\small
\begin{tabularx}{\linewidth}{@{}l X@{}}
\toprule
\textbf{File} & \textbf{Purpose} \\
\midrule
\code{SKILL.md} & The conversion procedure: ground rules, steps, writing conventions, budget. \\
\code{scripts/prepare.py} & Page strips, math detection, equation tags, printed page numbers. \\
\code{scripts/figures.py} & Figure detection, rendering and crop-check images. \\
\code{scripts/assemble.py} & Joins pages, adds the legend and page markers, packs the figures PDF. \\
\code{scripts/check.py} & KaTeX syntax, equation-tag, debris and page-marker checks. \\
\code{scripts/verify_eqs.py} & Side-by-side proofreading sheets for displayed equations. \\
\code{scripts/eval/mathcmp.py} & Regression test: formulas against a LaTeX source. \\
\code{scripts/eval/prosecmp.py} & Regression test: prose against a LaTeX source. \\
\bottomrule
\end{tabularx}
\end{center}

\end{document}
